\section*{Hoofdstuk \ref{ch:g6:solutions-and-solubility} --- Oplossingen en oplosbaarheid}

\begin{solution}{exo:g6:solutions-and-solubility:1}
\omterm{def:g6:solutions-and-solubility:solute}{Opgeloste stof}: suiker. \omterm{def:g6:solutions-and-solubility:solute}{Oplosmiddel}: water. Ja, het \omterm{def:g6:solutions-and-solubility:solute}{oplosmiddel} is
water: het is een \omterm{def:g6:solutions-and-solubility:solute}{waterige oplossing}.
\end{solution}

\begin{solution}{exo:g6:solutions-and-solubility:2}
Een oplossing die niets meer van haar \omterm{def:g6:solutions-and-solubility:solute}{opgeloste stof} kan \omterm{def:g2:mixing-and-dissolving:dissolve}{oplossen}. Wat
je er nog bij doet, blijft onopgelost op de bodem liggen.
\end{solution}

\begin{solution}{exo:g6:solutions-and-solubility:3}
$250 + 20 = \qty{270}{g}$.
\end{solution}

\begin{solution}{exo:g6:solutions-and-solubility:4}
Ethanol: \omterm{def:g6:solutions-and-solubility:miscible}{mengbaar}. Slaolie: \omterm{def:g6:solutions-and-solubility:miscible}{niet-mengbaar}. Azijn: \omterm{def:g6:solutions-and-solubility:miscible}{mengbaar} (het is
grotendeels water).
\end{solution}

\begin{solution}{exo:g6:solutions-and-solubility:5}
Propanon (aceton). GHS02: licht ontvlambaar; GHS07: irriteert de ogen,
de damp kan slaperigheid of duizeligheid veroorzaken.
\end{solution}

\begin{solution}{exo:g6:solutions-and-solubility:6}
$360 \times \frac{500}{1000} = \qty{180}{g}$.
\end{solution}

\begin{solution}{exo:g6:solutions-and-solubility:7}
\qty{200}{g} water lost hoogstens $360 \times \frac{200}{1000} =
\qty{72}{g}$ op. Nee: er blijft $80 - 72 = \qty{8}{g}$ op de bodem
liggen.
\end{solution}

\begin{solution}{exo:g6:solutions-and-solubility:8}
De oplossing weegt $225 + 25 = \qty{250}{g}$; $\frac{25}{250} =
\qty{10}{\percent}$.
\end{solution}

\begin{solution}{exo:g6:solutions-and-solubility:9}
In het eerste bekerglas, waar al het zout is opgelost en de oplossing nog
niet verzadigd is.
\end{solution}

\begin{solution}{exo:g6:solutions-and-solubility:10}
Het koolstofdioxide was onder druk in de gesloten fles opgelost. Als je
de fles opent, daalt de druk: het water kan minder gas vasthouden, en het
gas dat te veel is, komt eruit als belletjes.
\end{solution}

\begin{solution}{exo:g6:solutions-and-solubility:11}
$2000 \div 0.04 = 50\,000$ keer meer.
\end{solution}

\begin{solution}{exo:g6:solutions-and-solubility:12}
Een gas lost minder goed op in warm water: de warme vijver bevat minder
opgeloste zuurstof, en de vissen komen zuurstof tekort.
\end{solution}

\begin{solution}{pb:g6:solutions-and-solubility:1}
\textbf{1.} \omterm{def:g6:solutions-and-solubility:solute}{Opgeloste stof}: zout. \omterm{def:g6:solutions-and-solubility:solute}{Oplosmiddel}: water.

\textbf{2.} Het water bevat zoveel zout als het kan: er kan geen zout
meer in \omterm{def:g2:mixing-and-dissolving:dissolve}{oplossen}.

\textbf{3.} Het onopgeloste zout op de bodem bewijst dat de pekel
verzadigd is: ze bevat zoveel zout als mogelijk is.

\textbf{4.} $360 \times 20 = \qty{7200}{g} = \qty{7.2}{kg}$.

\textbf{5.} $20 + 7.2 = \qty{27.2}{kg}$.

\textbf{6.} $\frac{7.2}{27.2} \approx 0.26$: ongeveer \qty{26}{\percent}.

\textbf{7.} Nee: het water zou \qty{7.2}{kg} kunnen \omterm{def:g2:mixing-and-dissolving:dissolve}{oplossen}. Er kan nog
$7.2 - 5 = \qty{2.2}{kg}$ in \omterm{def:g2:mixing-and-dissolving:dissolve}{oplossen}.

\textbf{8.} $20 - 5 = \qty{15}{kg}$ water (\qty{15}{L}).

\textbf{9.} $360 \times 15 = \qty{5400}{g} = \qty{5.4}{kg}$.

\textbf{10.} Het komt als vaste kristallen uit de oplossing: het
kristalliseert uit en zakt naar de bodem.

\textbf{11.} Er verschijnen nieuwe witte kristallen, die zich op de bodem
van de bak ophopen.

\textbf{12.} Er is $7.2 - 5.4 = \qty{1.8}{kg}$ zout uitgekristalliseerd.
\end{solution}
